\newpage
\reviewsection{Future Classroom Incorporation of the Resources}

In my future classroom, I would use Kit for Kids during beginning-of-year community building and revisit the materials before coding, robotics, group projects, recess transitions, digital publishing, project showcases, and partner troubleshooting. These are moments when students must coordinate language, space, noise, materials, timing, and emotion. They are also moments when misunderstanding can become teasing, taking over, exclusion, or the assumption that another student is being difficult.

I would pair the OAR resources with supports available to the whole class: first/next/last directions, visual choices, quiet options, clear partner roles, help cards, predictable wait time, permission to pass, and multiple ways to respond. These supports should not appear only after one learner struggles or discloses a diagnosis. They should be part of the classroom's normal design.

\begin{figure}[htbp]
\centering
\resizebox{0.94\textwidth}{!}{\begin{tikzpicture}[stage/.style={rounded corners=3mm,draw=DeepBlue,very thick,minimum width=2in,minimum height=.85in,align=center},arrow/.style={-{Latex[length=3mm]},very thick,draw=PrimaryRed}]
\node[stage,fill=SoftBlue] (a) {\textbf{1. Prepare}\\visual routines, choices, quiet options};
\node[stage,fill=SoftGreen,right=7mm of a] (b) {\textbf{2. Teach}\\invite, wait, ask, respect, include};
\node[stage,fill=SoftGold,right=7mm of b] (c) {\textbf{3. Practice}\\coding, robotics, transitions, groups};
\node[stage,fill=SoftLavender,below=12mm of c] (d) {\textbf{4. Notice}\\communication, sensory load, pacing, belonging};
\node[stage,fill=SoftGray,left=7mm of d] (e) {\textbf{5. Adjust}\\environment, roles, directions, wait time};
\node[stage,fill=SoftBlue,left=7mm of e] (f) {\textbf{6. Repeat}\\make support a classroom routine};
\draw[arrow] (a)--(b); \draw[arrow] (b)--(c); \draw[arrow] (c)--(d); \draw[arrow] (d)--(e); \draw[arrow] (e)--(f); \draw[arrow] (f.west)..controls+(-1,0)and+(-1,0)..(a.west);
\end{tikzpicture}}
\caption{Future classroom integration cycle. OAR resources become useful when peer language is repeatedly taught, practiced, observed, and connected to environmental adjustments.}
\label{fig:future-cycle}
\end{figure}

Figure~\ref{fig:future-cycle} shows how I would prevent the materials from becoming a one-day awareness activity. The cycle begins with accessible classroom design, teaches concrete peer language, creates repeated opportunities to practice, and uses observation to adjust the environment. Repetition makes support part of classroom culture.

I would also make the five peer actions visible near collaborative work areas. Invitation must include consent, help should preserve ownership, waiting is an active form of respect, sensory tools are not toys, and different participation is still real participation. When developmentally appropriate, I would include autistic-authored or autistic-led perspectives so students do not learn about autism only through non-autistic interpretation.

\begin{center}
\colorbox{SoftGreen}{\parbox{0.92\textwidth}{\textbf{Future classroom commitment:} I will teach peer inclusion proactively, before conflict or exclusion occurs, and evaluate support by whether it expands communication, choice, competence, and belonging.}}
\end{center}

\reviewsection{Conclusion}

OAR's Kit for Kids materials provided a developmentally appropriate way to teach peer support inside ordinary classroom moments. Placement evidence helped make the scenario realistic and reinforced that behavior cannot be separated from environment, communication access, sensory conditions, and peer relationships.

My main takeaway is that positive behavioral support includes peer culture. A classroom should already know how to invite, wait, ask, respect, and include before harm happens. Autistic students and students with other complex support needs should not have to become more typical to be recognized as competent classmates. The classroom must become more responsive, curious, and flexible so different ways of communicating, moving, learning, and participating are treated as part of belonging.
